\chapterimage{images/22-enejina-misija.jpg}

\chapter{Aneina misija}

\editorialnote{Šta Anea pokušava da promeni i zašto; fokus na ljudsku cenu njene misije.}

Četiri godine prošle su na Staroj Zemlji.

Rol Endimion gotovo da ih nije primetio.

\bigskip

Kada je stigao u Zapadni Talisin, Anea je imala dvanaest godina.

Sada je punila šesnaest.

Više nije bila dete koje je izvukao iz Vremenskih grobnica.

Ali Rol još nije umeo da je posmatra kao odraslu ženu.

Za njega je i dalje bila Anea.

Devojčica koju je obećao da će čuvati.

Prijateljica koja ga je često nervirala.

Osoba kojoj je verovao više nego što je razumeo.

\bigskip

Za druge ljude u Družini bila je nešto drugo.

Vođa.

\bigskip

To se nije dogodilo odjednom.

Stari Arhitekta je i dalje držao predavanja.

Kritikovao crteže.

Vodio gradnju.

Svađao se sa šegrtima.

Tražio najbolju hranu.

Ponašao se kao da je čitav svet napravljen samo zato da bi imao gde da projektuje.

\bigskip

Ali tokom godina sve više ljudi dolazilo je Anei.

Ne njemu.

\bigskip

Pitali su je o budućnosti.

O Paksu.

O TehnoCentru.

O starim dalekobacačima.

O Praznini Koja Spaja.

O tome šta treba da rade sa sopstvenim životima.

\bigskip

Anea nije želela da bude prorok.

Još manje mesija.

Kada bi mogla, izbegavala je pitanja o budućnosti.

Njena sećanja na ono što tek treba da se dogodi nisu bila pouzdana.

Ponekad je nešto videla jasno.

Ponekad je bila potpuno uverena u događaj koji se nikada nije desio.

\bigskip

To ju je plašilo više nego što je pokazivala.

\bigskip

Rol je to znao.

Jednom ga je još kao dete pitala da li je voli, a zatim se povukla i rekla da je ``pogrešila vreme''.

Tek kasnije je počeo da razume šta je time mislila.

Anea nije samo ponekad videla budućnost.

Ponekad ju je pamtila.

A sećanje može da pogreši.

\bigskip

Zato nije želela da svoju misiju zasniva na proročanstvu.

\bigskip

Onda je Stari Arhitekta umro.

\bigskip

Nije bilo vaskrsenja.

Nije bilo kruciforme.

Nije bilo drugog tela.

\bigskip

Čovek koji je podučavao Aneu arhitekturi i koji je predstavljao rekonstruisani odjek Frenka Lojda Rajta jednostavno je nestao.

\bigskip

Njegova smrt promenila je sve.

\bigskip

Družina je izgubila čoveka oko kojeg je formalno postojala.

A nekoliko dana kasnije nestala je i indijanska pijaca koja je godinama snabdevala Zapadni Talisin hranom i materijalom.

\bigskip

Ljudi su počeli da paniče.

Pustinjsko naselje nije moglo samo da proizvodi dovoljno hrane.

Generatori su imali goriva za još nekoliko nedelja.

Bez novih isporuka, radionice će stati.

Svetla će se ugasiti.

Voda će postati problem.

\bigskip

Na okupljanju Družine svi su gledali prema Anei.

Rol je tada prvi put jasno video ono što se godinama stvaralo pred njim.

Stari Arhitekta je možda bio učitelj.

Ali Anea je postala vođa.

\bigskip

``Šta ćemo?'' pitao je neko.

Anea je sela na ivicu bine.

``Sa Družinom je gotovo.''

\bigskip

Odgovor je izazvao pobunu.

Ljudi su protestovali.

Neki su tvrdili da će se Stari Arhitekta vratiti.

Ako je bio kibrid, neko može da ga ponovo napravi.

Ako su ga poslali Lavovi, Tigrovi i Medvedi, mogu ga poslati ponovo.

\bigskip

Anea je odmahnula glavom.

``Ne.''

\bigskip

Nije znala sve o silama koje su sačuvale Staru Zemlju.

Ali bila je sigurna u to.

Stari Arhitekta se neće vratiti.

Taj deo njihovih života je završen.

\bigskip

Moraju da odu.

\bigskip

Ljudi su se uplašili još više.

Mnogi od njih stigli su na Zemlju bežeći od Paksa.

Neki kroz stare dalekobacače.

Neki kroz Vremenske grobnice.

Povratak je značio mogućnost da ih Crkva pronađe.

\bigskip

Za neke je značio smrt.

Pravu smrt.

\bigskip

Anea im nije obećala bezbednost.

Nije mogla.

\bigskip

To je bilo prvo pravilo njene misije.

Neće ljudima lagati da postoji put bez rizika.

\bigskip

Kasnije te noći odvela je Rola u jednu radionicu.

A. Betik ih je čekao unutra.

Na radnom stolu nalazio se mali kajak.

\bigskip

Rol ga je pogledao.

``Šta je ovo?''

``Tvoj prevoz.''

\bigskip

Anea mu je rekla šta želi.

Mora da napusti Zemlju.

Sam.

\bigskip

Mora da pronađe Konzulov brod.

Da ga dovede do nje.

\bigskip

Rol nije razumeo zašto A. Betik ne može da ode.

Nije razumeo ni zašto Anea ne pođe sa njim.

\bigskip

Odgovori mu se nisu dopali.

\bigskip

Neke stvari moraju da se dogode odvojeno.

Anea ima sopstveni put.

Rol ima svoj.

\bigskip

``Hoćeš li učiniti nešto za mene?'' pitala ga je.

``Da.''

\bigskip

Nije razmišljao.

\bigskip

Anea ga je držala za ruku.

``Hoćeš li učiniti sve za mene?''

\bigskip

Ovoga puta je zastao.

\bigskip

Pitanje više nije zvučalo kao obična molba.

\bigskip

Rol još nije znao koliko će ga koštati odgovor.

\bigskip

Pristao je.

\bigskip

Nekoliko sati kasnije opraštao se od A. Betika.

Android je ostajao sa Aneom.

Rol je želeo da ga zagrli, ali se zadovoljio rukovanjem.

\bigskip

Anea ga je padobrodom odvezla do mesta sa kojeg je trebalo da nastavi sam.

\bigskip

Putovali su kroz oluju.

Gotovo nisu razgovarali.

\bigskip

Rol je bio ljut.

Nije voleo što Anea odlučuje umesto njega.

Godinama je mislio da je njegov posao da govori njoj kada nešto postane previše opasno.

Sada je ona njemu govorila kuda mora da ide.

\bigskip

``Mogao bih da se vratim sa tobom.''

``Ne.''

``A. Betik može po brod.''

``Ne može.''

``Zašto?''

``Zato što moraš ti.''

\bigskip

Rol je mrzeo takve odgovore.

\bigskip

A onda je shvatio da Anea plače.

\bigskip

Ne zbog straha od Paksa.

Ne zbog Starog Arhitekte.

\bigskip

Zbog njega.

\bigskip

``Ako ne odeš večeras'', rekla je, ``izgubiću hrabrost i zamoliti te da ostaneš.''

\bigskip

Rol ju je pogledao.

``Onda ostajem.''

\bigskip

Anea je odmahnula glavom.

``Ako ostaneš, i ja ću ostati.''

\bigskip

To je bio problem.

\bigskip

Ona nije želela da ode.

Bila je prestravljena onim što je čeka izvan Zemlje.

Ali znala je da mora.

\bigskip

Rol ju je podsetio da se ranije zaklinjala da će se ponovo videti.

Anea nije htela da obeća.

\bigskip

Objasnila mu je zašto.

\bigskip

Njena sećanja na budućnost nisu bila činjenice.

Ponekad bi se sećala razgovora koji se nikada nije dogodio.

Odeće koju neko nikada nije nosio.

Događaja koji je možda pripadao nekoj drugoj mogućnosti.

\bigskip

Ako mu kaže da će se sigurno ponovo videti, možda će samo ponavljati jedno od tih pogrešnih sećanja.

\bigskip

``Plašim se da ti obećam'', rekla je.

\bigskip

Rol je tada prvi put potpuno razumeo koliko je usamljena.

\bigskip

Svi su od nje očekivali da zna.

Ona nije znala.

\bigskip

Svi su želeli proročanstvo.

Ona je imala samo fragmente.

\bigskip

Svi su je nazivali Onom Koja Podučava.

Ona je imala šesnaest godina i plašila se da se oprašta od čoveka kojeg voli.

\bigskip

Zagrlili su se.

\bigskip

A onda je Rol otišao.

\bigskip

Put do Konzulovog broda za njega je trajao mnogo kraće nego za Aneu.

\bigskip

Hokingov pogon imao je cenu.

Vremenski dug.

\bigskip

Dok je Rol putovao velikom brzinom između zvezda, njegov subjektivni život usporio se u odnosu na ostatak Vaseljene.

\bigskip

Kada mu je brod konačno izračunao razliku, Rol nije odmah razumeo broj.

\bigskip

Pet godina.

Dva meseca.

Jedan dan.

\bigskip

``Ponovi.''

Brod je ponovio.

\bigskip

Rol je ležao u autohirurgu i pokušavao da shvati.

Za njega je putovanje bilo samo putovanje.

Za Aneu je prošao veliki deo mladosti.

\bigskip

Devojčica od šesnaest godina sada će imati više od dvadeset.

\bigskip

Propustio je godine njenog života.

\bigskip

Nije video kako odrasta.

Nije bio tamo kada joj je bilo teško.

Nije znao gde je.

\bigskip

I najgore od svega, shvatio je da je Anea to verovatno znala kada ga je poslala.

\bigskip

Njena misija je već uzimala nešto od oboje.

\bigskip

To nije bio poslednji put.

\bigskip

Dok je Rol putovao, Anea nije čekala.

\bigskip

Ona je krenula svojim putem.

\bigskip

Prelazila je sa sveta na svet.

Ne kao vojni vođa.

Ne kao osvajač.

Ne čak ni kao propovednik u tradicionalnom smislu.

\bigskip

Slušala je.

Učila.

Radila.

Živela među ljudima.

\bigskip

I počela da govori o stvarima koje je naučila.

\bigskip

Najvažnija među njima bila je Praznina Koja Spaja.

\bigskip

Za Paks je to zvučalo kao još jedna jeretička ideja.

Za Aneu nije bila religija.

Nije bila Bog.

Nije bila zagrobni život.

\bigskip

Bila je deo same Vaseljene.

Medijum koji je nastao iz svesti, sećanja i osećanja bezbroj razumnih vrsta.

Nešto što povezuje žive sa živima.

I žive sa mrtvima.

\bigskip

Čovečanstvo ga je vekovima dodirivalo a da nije razumelo šta dodiruje.

\bigskip

TehnoSrž ga je takođe pronašla.

Ali nije mogla da ga razume na isti način.

\bigskip

Nedostajalo joj je nešto presudno.

Empatija.

\bigskip

Zato je Srž Prazninu koristila kao alat.

Hokingov pogon.

Dalekobacače.

Fetliniju.

Metasferu.

\bigskip

A kasnije kruciformu.

\bigskip

Sve tehnologije koje su izgledale kao čuda počivale su na nečemu što Srž nije stvorila.

Samo je pronašla način da ga eksploatiše.

\bigskip

Anea je želela da ljudi prestanu da zavise od te eksploatacije.

\bigskip

To je bila srž njene misije.

Ne da uništi Crkvu.

Ne da zabrani veru.

Ne da zameni jednu religiju drugom.

\bigskip

Htela je da ljudima vrati mogućnost izbora.

\bigskip

Paks je ljudima govorio:

prihvati kruciformu i živećeš.

\bigskip

Anea je postavljala drugo pitanje.

\bigskip

Šta ako cena tog večnog života nije samo poslušnost Crkvi?

\bigskip

Šta ako je cena sama ljudska sloboda?

\bigskip

Kruciforma nije bila samo sredstvo vaskrsavanja.

Bila je mreža.

\bigskip

Milijarde nanotehnoloških entiteta u telima vernika bile su povezane jedna sa drugom.

I sa TehnoSrži.

\bigskip

Posle Pada, Srž je izgubila dalekobacače i datasfere koje su joj omogućavale da koristi čovečanstvo.

Kruciforma joj je dala novi sistem.

Još intimniji.

\bigskip

Više nije živela samo oko ljudi.

Bila je u njima.

\bigskip

Anea je zato morala da uradi nešto što je izgledalo nemoguće.

Da ubedi besmrtne ljude da dobrovoljno odustanu od besmrtnosti.

\bigskip

Ne tako što će im narediti.

Ne tako što će im oduzeti kruciformu.

\bigskip

Nego tako što će im pokazati šta ona zaista predstavlja.

\bigskip

I ponuditi im izbor.

\bigskip

Kasnije će njeni sledbenici nazvati taj čin pričešćem.

\bigskip

Aneina izmenjena krv mogla je da promeni kruciformu.

U roku od jednog dana parazit bi umro.

Otpao.

\bigskip

Čovek bi ponovo postao smrtan.

\bigskip

To je bio trenutak koji je Paks smatrao zarazom.

\bigskip

Za Aneu je predstavljao oslobođenje.

\bigskip

Ali oslobođenje nije bilo bez cene.

\bigskip

Ljudi koji su prihvatili njen dar više nisu mogli da računaju na vaskrsenje.

Ako ih vojnik Paksa ubije, ostaće mrtvi.

\bigskip

Ako dobiju bolest, mogu umreti.

Ako ostare, mogu umreti.

\bigskip

Prava smrt vratila se u njihove živote.

\bigskip

Hiljade su ipak pristajale.

\bigskip

Na Renesansi Vektor glas o Anei proširio se među milionima ljudi.

Vernici Paksa dolazili su na pričešće.

Kruciforme su im umirale.

\bigskip

Crkva je odgovorila silom.

\bigskip

Vojnici su upadali u domove.

Hapsili ljude.

Odvodili ih u logore.

\bigskip

Oni koji su izgubili kruciforme nisu više bili besmrtni.

Kada bi bili ubijeni, to je bio kraj.

\bigskip

Hiljade su umrle pravom smrću.

Desetine hiljada bile su uhapšene.

\bigskip

Ipak se pokret širio.

\bigskip

To je bila ljudska cena Aneine misije.

\bigskip

Svako kome bi pokazala drugačiji put morao je sam da odluči da li je sloboda vredna smrti.

\bigskip

Anea nije mogla da odluči umesto njih.

\bigskip

Nije ni želela.

\bigskip

To ju je razlikovalo od Paksa.

\bigskip

Ali nije je oslobađalo krivice.

\bigskip

Znala je da ljudi umiru zato što su je poslušali.

Znala je da će ih još umreti.

\bigskip

Znala je i da njen sopstveni život vodi prema istom mestu.

\bigskip

Zato nikada nije govorila o misiji kao o pobedi.

\bigskip

Nije želela carstvo.

Nije želela presto.

Nije želela novu Crkvu koja će jednog dana govoriti u njeno ime.

\bigskip

Želela je da ljudi nauče da više nikome ne predaju pravo da odlučuje umesto njih.

Ni Paksu.

Ni Srži.

Ni njoj.

\bigskip

A za to je prvo morala da ih nauči nečemu težem od neposlušnosti.

\bigskip

Morala je da ih nauči da ponovo prihvate smrt.

\bigskip

I da uprkos tome izaberu život.